% มคอ.5 — 7015906 ภาคเรียนที่ 2/2568
\documentclass{mko5}
\begin{document}
\mkofivetitle

\coursecode{7015906}
\coursename{ระเบียบวิธีวิจัยทางวิทยาศาสตร์และวิศวกรรมศาสตร์}
\curriculumnote{หลักสูตรวิศวกรรมศาสตรมหาบัณฑิต สาขาวิชาวิศวกรรมคอมพิวเตอร์และปัญญาประดิษฐ์ พ.ศ. 2568}
\coursegroup{วิชาเฉพาะด้านบังคับ}
\coursecredits{3 (2-2-5)}
\studyplan{แผน 2 --- วิชาชีพ (สารนิพนธ์)}
\semester{ภาคเรียนที่ 2 / ปีการศึกษา 2568}
\prerequisite{ไม่มี}
\corequisite{ไม่มี}
\studyplace{มหาวิทยาลัยราชภัฏอุตรดิตถ์ --- คณะเทคโนโลยีอุตสาหกรรม}
\registeredstudents{4 คน}
\completedstudents{4 คน}
\instructor{ผศ.ดร.กาญจนา ดาวเด่น}
\coinstructors{---}

\mkoheaderinfo

\begin{teachingactualtable}
  \actualrow{CLO1}{บรรยายหลักการระเบียบวิธีวิจัย อภิปรายกลุ่ม แบบทดสอบและแบบฝึกหัด สัปดาห์ 1--2}
  \actualrow{CLO2}{สาธิตการค้นคว้าวิจัย Workshop ทบทวนวรรณกรรม ตั้งสมมติฐ์ ออกแบบการวิจัย โครงการวิจัยสัปดาห์ 11--15}
  \actualrow{CLO3}{แนะนำการจัดการข้อมูล ฝึกเก็บรวบรวมข้อมูล รายงานจัดระบบข้อมูล นักศึกษาทุกคนส่งโครงร่างทบทวนวรรณกรรมแล้ว}
  \actualrow{CLO4}{บรรยายจรรยาบรรณและจริยธรรม กรณีศึกษา บทความสะท้อนคิด การประเมินการอ้างอิง --- พัฒนาการ CLO4 ดี (สอดคล้อง PLO4)}
  \actualrow{CLO5}{Lab สถิติเชิงพรรณนาและเชิงอนุมาน ฝึกใช้ซอฟต์แวร์สถิติ รายงานผลการวิเคราะห์ข้อมูล}
\end{teachingactualtable}

\begin{assessmentusedtable}
  \assessusedrow{CLO1}{แบบทดสอบข้อเขียน การตอบคำถามในชั้นเรียน --- Rubric 4 ระดับ}
  \assessusedrow{CLO2}{รายงานทบทวนวรรณกรรม การนำเสนอแนวคิดงานวิจัย โครงการวิจัย/สอบปลายภาค}
  \assessusedrow{CLO3}{รายงานการจัดระบบข้อมูล การนำเสนอ รายงานความก้าวหน้าโครงการ}
  \assessusedrow{CLO4}{รายงานกรณีศึกษาด้านจริยธรรม การประเมินการเขียนอ้างอิง บทสะท้อนคิด}
  \assessusedrow{CLO5}{รายงานผลการวิเคราะห์ข้อมูล การนำเสนอผลการทดลอง สอบปลายภาค}
\end{assessmentusedtable}

\mkosectionthree

\closummaryblock{CLO1}{อธิบาย หลักการและระเบียบวิธีวิจัยทางวิทยาศาสตร์และวิศวกรรมศาสตร์ ได้อย่างถูกต้องตามหลักวิชาการ}{80}
\achievementpass{4}{100}{ผ่านเกณฑ์ Rubric ระดับ 3 ขึ้นไป}
\achievementfail{0}{0}{ไม่มี}
\closummaryend

\closummaryblock{CLO2}{วิเคราะห์ ปัญหาวิจัยและตั้งสมมติฐานการวิจัย ที่สอดคล้องกับบริบทของงานวิจัยทางวิศวกรรมคอมพิวเตอร์}{80}
\achievementpass{4}{100}{โครงร่างทบทวนวรรณกรรมและข้อเสนอโครงการผ่านเกณฑ์}
\achievementfail{0}{0}{ไม่มี}
\closummaryend

\closummaryblock{CLO3}{จัดระบบ ข้อมูลจากการทบทวนวรรณกรรมและการเก็บรวบรวมข้อมูล เพื่อนำไปสู่การพัฒนาโจทย์วิจัย}{80}
\achievementpass{4}{100}{รายงานจัดระบบข้อมูลและความก้าวหน้าโครงการผ่าน}
\achievementfail{0}{0}{ไม่มี}
\closummaryend

\closummaryblock{CLO4}{แสดงออก ถึงการมีจริยธรรมในการทำวิจัยและการเขียนรายงาน ผ่านการอ้างอิงผลงานทางวิชาการอย่างถูกต้อง}{80}
\achievementpass{4}{100}{ผ่านเกณฑ์ CLO4 (Affective 5) ทุกคน}
\achievementfail{0}{0}{ไม่มี}
\closummaryend

\closummaryblock{CLO5}{ประยุกต์ใช้ วิธีการทางสถิติในการวิเคราะห์ข้อมูลและสรุปผลการวิจัย ได้อย่างเหมาะสมกับประเภทของข้อมูล}{80}
\achievementpass{4}{100}{รายงานวิเคราะห์ข้อมูลและสอบปลายภาคผ่าน}
\achievementfail{0}{0}{ไม่มี}
\closummaryend

\mkosectionfour
\begin{mkotextblock}
นักศึกษารุ่นแรก ($n=4$) ยังอยู่ระหว่างกำหนดหัวข้อวิจัย/สารนิพนธ์ขั้นสุดท้าย;
การใช้ซอฟต์แวร์สถิติบางกลุ่มต้องการเวลาฝึกเพิ่มเติมนอกชั้นเรียน
\end{mkotextblock}

\mkosectionfive
\begin{mkotextblock}
จัดทำตัวอย่างโครงร่างทบทวนวรรณกรรมฉบับมาตรฐาน;
เพิ่ม Workshop การใช้เครื่องมือสถิติ (Python/R) ก่อนสอบปลายภาค;
เชื่อมโยงผล มคอ.5 กับการกำหนดหัวข้อวิจัยของนักศึกษาแต่ละรายในภาคเรียนที่ 3/2568
\end{mkotextblock}

\mkosectionsix
\begin{mkotextblock}
ขอให้ห้องสมุด/ฐานข้อมูลวิจัยสนับสนุนการเข้าถึง Scopus/IEEE อย่างต่อเนื่อง;
พิจารณาจัดอบรมการเขียนอ้างอิง APA/IEEE สำหรับนักศึกษารุ่นแรก
\end{mkotextblock}

\mkosectionseven
\mkosignaturepair{ผศ.ดร.กาญจนา ดาวเด่น}{30 เมษายน 2569}{ผศ.ดร.พิทักษ์ คล้ายชม}{30 เมษายน 2569}
\end{document}
